% YIN-VERBATIM-SECTION: video=OY_napMPywE; range=00:45:00.000--00:55:00.000
% YIN-VERBATIM-CANONICAL-TRANSCRIPT-SHA256: 5ac8ac5fb25a3235d8fa11b2b6be99b5f2bb9329d307c4045629544f4e43e9bd
% YIN-VERBATIM-NOTES: 253a:4--5; pdf=9--10

% YIN-VERBATIM: id=YIN-OY-T000166; time=00:45:00.000--00:45:01.430; speaker=NONE; retained_words=0; disposition=rolling_caption_tail_no_new_speech
% YIN-VERBATIM: id=YIN-OY-T000167; time=00:45:01.430--00:45:01.440; speaker=NONE; retained_words=0; disposition=rolling_caption_separator_no_new_tokens
% YIN-VERBATIM: id=YIN-OY-T000168; time=00:45:01.440--00:45:07.490; speaker=XI_YIN; retained_words=4; disposition=classroom_question_pause
\noindent\textsc{Xi Yin:} Okay. Any other questions?\par

% YIN-VERBATIM: id=YIN-OY-T000169; time=00:45:07.490--00:45:07.500; speaker=NONE; retained_words=0; disposition=rolling_caption_separator_no_new_tokens
% YIN-VERBATIM: id=YIN-OY-T000170; time=00:45:07.500--00:45:31.800; speaker=XI_YIN; retained_words=30; disposition=included_with_uncertain_clause
% YIN-UNCERTAINTY: Exact wording of the Hilbert-space/operator clause is unresolved; the bracket records source-supported content without claiming verbatim speech.
\noindent\textsc{Xi Yin:} All right. So here it is. We have [unclear: a well-defined system, its operators, Hilbert space, and Hamiltonian]. And by design, this is a quantum-mechanical system of free relativistic particles.\par

% YIN-VERBATIM: id=YIN-OY-T000171; time=00:45:31.800--00:45:31.810; speaker=NONE; retained_words=0; disposition=rolling_caption_separator_no_new_tokens
% YIN-VERBATIM: id=YIN-OY-T000172; time=00:45:31.810--00:45:34.430; speaker=NONE; retained_words=0; disposition=spurious_nonspeech_caption
% YIN-VERBATIM: id=YIN-OY-T000173; time=00:45:34.430--00:45:34.440; speaker=NONE; retained_words=0; disposition=rolling_caption_separator_no_new_tokens
% YIN-VERBATIM: id=YIN-OY-T000174; time=00:45:34.440--00:46:08.030; speaker=XI_YIN; retained_words=59; disposition=included
\noindent\textsc{Xi Yin:} So if you're just interested in describing quantum mechanics of free particles, free in the sense that they're not interacting with one another, but then we're done. And so then you don't have to take this course; you can go home. But the thing is, we would like to be able to describe interactions. So what about interacting particles?\par

% YIN-VERBATIM: id=YIN-OY-T000175; time=00:46:08.030--00:46:08.040; speaker=NONE; retained_words=0; disposition=rolling_caption_separator_no_new_tokens
% YIN-VERBATIM: id=YIN-OY-T000176; time=00:46:08.040--00:46:31.370; speaker=XI_YIN; retained_words=50; disposition=included
% YIN-UNCERTAINTY: The first spoken name for the free Hamiltonian sounds like H-free but is less clear than the subsequent H-zero.
\noindent\textsc{Xi Yin:} Well, you might say that that's not a big deal. So, you know, I can take, so let's call this $H_{\mathrm{free}}$; let's call this $H_0$, for the free theory. So in quantum mechanics we can just pick whatever Hermitian operator we would like, so you can add some interaction terms.\par

% YIN-EXACT-NOTE-EQUATION-PLACEHOLDER: notes=253a:4; pdf=9; content=H=H_0+H_{\mathrm{int}}.
% YIN-EXACT-NOTE-EQUATION-PLACEHOLDER: notes=253a:4; pdf=9; source-content=a^+a^+a,\;aaa^+,\;\ldots; transcript-normalization=\dagger.

% YIN-VERBATIM: id=YIN-OY-T000177; time=00:46:31.370--00:46:31.380; speaker=NONE; retained_words=0; disposition=rolling_caption_separator_no_new_tokens
% YIN-VERBATIM: id=YIN-OY-T000178; time=00:46:31.380--00:47:15.170; speaker=XI_YIN; retained_words=90; disposition=included_with_spoken_self_correction
% YIN-UNCERTAINTY: The captions do not preserve the cubic monomials. Their exact order follows handwritten 253a:4; the transcript normalizes the creator mark to \dagger.
\noindent\textsc{Xi Yin:} And let's try that. Okay, so, for example, you might even include interaction terms in the Hamiltonian that allow changing particle numbers. So, for example, you can have operators like $a^\dagger a^\dagger a$, with some momentum assignments. Sorry, let's say like this, or $a a a^\dagger$. So this guy, for example, can take this one-particle state, annihilate the one particle, and create two particles. So this operator would act on the one-particle state and give you a two-particle state. This internal kind of interaction will change the number of particles.\par

% YIN-VERBATIM: id=YIN-OY-T000179; time=00:47:15.170--00:47:15.180; speaker=NONE; retained_words=0; disposition=rolling_caption_separator_no_new_tokens
% YIN-VERBATIM: id=YIN-OY-T000180; time=00:47:15.180--00:47:23.030; speaker=XI_YIN; retained_words=11; disposition=included_transition
\noindent\textsc{Xi Yin:} Okay. So we can do that. So what's the big deal?\par

% YIN-VERBATIM: id=YIN-OY-T000181; time=00:47:23.030--00:47:23.040; speaker=NONE; retained_words=0; disposition=rolling_caption_separator_no_new_tokens
% YIN-VERBATIM: id=YIN-OY-T000182; time=00:47:23.040--00:48:05.390; speaker=XI_YIN; retained_words=51; disposition=included
\noindent\textsc{Xi Yin:} The issue is the following. So in principle you could have done that, but it's quite difficult to introduce interactions in this manner in a way that preserves relativistic symmetry. Okay, so the issue is that it is possible, but not easy, not easy to find $H_{\mathrm{int}}$ that respects relativistic symmetry.\par

% YIN-VERBATIM: id=YIN-OY-T000183; time=00:48:05.390--00:48:05.400; speaker=NONE; retained_words=0; disposition=rolling_caption_separator_no_new_tokens
% YIN-VERBATIM: id=YIN-OY-T000184; time=00:48:05.400--00:48:21.000; speaker=XI_YIN; retained_words=13; disposition=included_transition_with_uncertain_word
% YIN-UNCERTAINTY: One adjective at 00:48:12.900--00:48:16.250 remains unresolved and has no mathematical consequence.
\noindent\textsc{Xi Yin:} Why is that? Why do I say that? The issue is quite [unclear].\par

% YIN-EXACT-NOTE-FIGURE-PLACEHOLDER: notes=253a:5; pdf=10; kind=spacetime-lightcone.
% YIN-EXACT-NOTE-FIGURE-REQUIREMENTS: time-axis=x^0; space-axis=\vec x; lightcone; timelike ray; spacelike red ray; label=super-luminal propagation forbidden.

% YIN-VERBATIM: id=YIN-OY-T000185; time=00:48:21.000--00:49:24.670; speaker=XI_YIN; retained_words=98; disposition=included
\noindent\textsc{Xi Yin:} Okay, so let me draw a schematic picture of space and time, space labeled by some spatial coordinate $\vec x$, and time. So if we have some event or some disturbance of the system that occurs at one point in space, say at time zero at some location in space, let's say here's the light cone. A light ray travels in this picture. I'm only drawing one spatial direction, right? So I'm also working in units where the speed of light is equal to one, so light will travel along rays at forty-five-degree angles, left and right, and...\par

% YIN-VERBATIM: id=YIN-OY-T000186; time=00:49:24.670--00:49:24.680; speaker=NONE; retained_words=0; disposition=rolling_caption_separator_no_new_tokens
% YIN-VERBATIM: id=YIN-OY-T000187; time=00:49:24.680--00:49:57.890; speaker=XI_YIN; retained_words=51; disposition=included
\noindent\textsc{Xi Yin:} A trajectory like this would be timelike; it would correspond to propagation of a signal slower than the speed of light. Whereas a trajectory like this, outside of the light cone, would be superluminal, and that's not supposed to be allowed in a relativistic theory. So superluminal propagation should be forbidden.\par

% YIN-VERBATIM: id=YIN-OY-T000188; time=00:49:57.890--00:49:57.900; speaker=NONE; retained_words=0; disposition=rolling_caption_separator_no_new_tokens
% YIN-VERBATIM: id=YIN-OY-T000189; time=00:49:57.900--00:50:23.890; speaker=XI_YIN; retained_words=64; disposition=included
\noindent\textsc{Xi Yin:} So I'm not allowed to say that I can, you know, disturb the system here, there's a red point here, and measure some response at that point in spacetime. That's not possible. If I disturb the system over here, I should only have a response within the light cone. That's just a constraint from relativistic causality, and any relativistic theory should have this property.\par

% YIN-VERBATIM: id=YIN-OY-T000190; time=00:50:23.890--00:50:23.900; speaker=NONE; retained_words=0; disposition=rolling_caption_separator_no_new_tokens
% YIN-VERBATIM: id=YIN-OY-T000191; time=00:50:23.900--00:51:05.270; speaker=XI_YIN; retained_words=79; disposition=included
% YIN-UNCERTAINTY: The phrase a priori instantaneous is a phonetic repair of severely corrupted ASR, supported by the sentence and note-page statement that propagation can be instantaneous.
\noindent\textsc{Xi Yin:} But you see, this kind of property is not built in in the form of the Hamiltonian formulation. Nothing tells you the speed limit. You know, the influence will be a priori instantaneous. In fact, this is not a very easy exercise, but you can try; you're welcome to. Try to include some interaction terms and see if you can arrange them such that there will be no superluminal propagation of signals. And that's actually not an easy thing.\par

% YIN-VERBATIM: id=YIN-OY-T000192; time=00:51:05.270--00:51:05.280; speaker=NONE; retained_words=0; disposition=rolling_caption_separator_no_new_tokens
% YIN-VERBATIM: id=YIN-OY-T000193; time=00:51:05.280--00:51:19.309; speaker=XI_YIN; retained_words=16; disposition=included_transition
\noindent\textsc{Xi Yin:} But let me quantify this a bit better in the language of these local field operators.\par

% YIN-VERBATIM: id=YIN-OY-T000194; time=00:51:19.309--00:51:19.319; speaker=NONE; retained_words=0; disposition=rolling_caption_separator_no_new_tokens
% YIN-VERBATIM: id=YIN-OY-T000195; time=00:51:19.319--00:52:17.630; speaker=XI_YIN; retained_words=74; disposition=included
\noindent\textsc{Xi Yin:} So in order to actually describe this kind of a local disturbance of the system in a way that is mathematically meaningful, so in a relativistic quantum-mechanical system we might expect the local disturbance... And so you change the state of your system by some operation localized at a point in space, at an instant in time. You might expect that there's a local disturbance to be represented by some kind of field operator.\par

% YIN-VERBATIM: id=YIN-OY-T000196; time=00:52:17.630--00:52:17.640; speaker=NONE; retained_words=0; disposition=rolling_caption_separator_no_new_tokens
% YIN-VERBATIM: id=YIN-OY-T000197; time=00:52:17.640--00:52:48.109; speaker=XI_YIN; retained_words=75; disposition=included
\noindent\textsc{Xi Yin:} And then when I say the word field, what does it mean? All I mean is that the operator depends on some spacetime point $x$. So that is to say, we have some operator. This operator, in the usual sense of quantum mechanics, acts on the Hilbert space, and we can take it to be Hermitian, if you wish. And, you know, its expectation value is supposed to be some physical observable, in some sense.\par

% YIN-VERBATIM: id=YIN-OY-T000198; time=00:52:48.109--00:52:48.119; speaker=NONE; retained_words=0; disposition=rolling_caption_separator_no_new_tokens
% YIN-VERBATIM: id=YIN-OY-T000199; time=00:52:48.119--00:53:15.000; speaker=XI_YIN; retained_words=64; disposition=included
% YIN-UNCERTAINTY: The opening noun operators is reconstructed from syntax; the remainder of the distinction is clear.
\noindent\textsc{Xi Yin:} But unlike the operators perhaps you are familiar with, what we have here is actually an entire family of operators labeled by some parameter. So I should emphasize that in my notation here. Unlike probably your undergraduate quantum mechanics course, where $x$ and $p$ are operators, here when I write $x$ in this context, they're just parameters; they're not operators. The operator is $\hat\phi$.\par

% YIN-EXACT-NOTE-EQUATION-PLACEHOLDER: notes=253a:5; pdf=10; content=x^\mu=(x^0,\vec x); qualification=x^\mu are parameters, not operators.

% YIN-VERBATIM: id=YIN-OY-T000200; time=00:53:15.000--00:53:39.470; speaker=XI_YIN; retained_words=57; disposition=included
\noindent\textsc{Xi Yin:} All right. So this $x^\mu$, which is some spacetime point, specified by the time $x^0$ and the spatial coordinate $\vec x$, this is just a label. This is just a label, a spacetime point. And I want to say that associated with that spacetime point there is some operator $\hat\phi$ that represents a disturbance of the system.\par

% YIN-VERBATIM: id=YIN-OY-T000201; time=00:53:39.470--00:53:39.480; speaker=NONE; retained_words=0; disposition=rolling_caption_separator_no_new_tokens
% YIN-VERBATIM: id=YIN-OY-T000202; time=00:53:39.480--00:53:52.010; speaker=XI_YIN; retained_words=4; disposition=classroom_question_pause
\noindent\textsc{Xi Yin:} Any questions so far?\par

% YIN-VERBATIM: id=YIN-OY-T000203; time=00:53:52.010--00:53:52.020; speaker=NONE; retained_words=0; disposition=rolling_caption_separator_no_new_tokens
% YIN-VERBATIM: id=YIN-OY-T000204; time=00:53:52.020--00:54:06.770; speaker=XI_YIN; retained_words=26; disposition=included_transition
\noindent\textsc{Xi Yin:} Now this field operator should obey some properties. So first of all, I said this in words early on, but let me say this more precisely.\par

% YIN-VERBATIM: id=YIN-OY-T000205; time=00:54:06.770--00:54:06.780; speaker=NONE; retained_words=0; disposition=rolling_caption_separator_no_new_tokens
% YIN-VERBATIM: id=YIN-OY-T000206; time=00:54:06.780--00:54:34.849; speaker=XI_YIN; retained_words=30; disposition=included
\noindent\textsc{Xi Yin:} So if I look at $\hat\phi(x)$ and $\hat\phi(x')$, at the different point $x'$, these two operators are supposed to be related by a Poincaré symmetry that sends $x$ to $x'$.\par

% YIN-VERBATIM: id=YIN-OY-T000207; time=00:54:34.849--00:54:34.859; speaker=NONE; retained_words=0; disposition=rolling_caption_separator_no_new_tokens
% YIN-VERBATIM: id=YIN-OY-T000208; time=00:54:34.859--00:55:00.000; speaker=XI_YIN; retained_words=61; disposition=included_continues
% YIN-UNCERTAINTY: The interval ends after acting. No covariance equation or infinitesimal-generator signs occur before this cutoff.
\noindent\textsc{Xi Yin:} So, you know, given two spacetime points, and because we're in Minkowski spacetime, I can always find some Poincaré symmetry. This could be just a translation; it could be a different Lorentz transformation, whatever. Of course it's highly nonunique. That takes $x$ to $x'$. And you would expect, if you have a symmetry that takes $x$ to $x'$, this symmetry, acting...\par

% YIN-NEXT-NOTE-EQUATION-OUTSIDE-LANE: notes=253a:6; pdf=11; the explicit covariance equation begins after 00:55:00.000.
% YIN-VERBATIM-WORD-COUNT: eligible=1007; retained=1007; ratio=1.000
